% v2_ch3_pipeline.tex — Chapter 3: Pipeline Architecture
% Video-to-Gaussian: An Agentic Pipeline for Volumetric Cultural Heritage Preservation
% Dr John O'Hare, DreamLab AI Consulting Ltd — June 2026

\chapter{Pipeline Architecture}\label{ch:pipeline}

The \texttt{gaussian-toolkit} is not a shell script that glues together pre-existing tools;
it is a purposefully architected system in which an \acrshort{llm} orchestrator reasons
about state, selects tool calls, handles retries, and passes structured artefacts between
bounded contexts. This chapter documents the design principles that shaped that architecture,
the eight processing stages executed for every acquisition, the three-container compute
fabric that hosts them, the hardware platform, and the control surfaces through which
operators interact with the running system.

% ─────────────────────────────────────────────────────────────────────────────
\section{Design Principles}\label{sec:design-principles}
% ─────────────────────────────────────────────────────────────────────────────

\subsection{Domain-Driven Design and Bounded Contexts}

The architecture follows \acrfull{ddd}, partitioning responsibility into seven
bounded contexts whose interfaces are explicit and whose internal implementations
may evolve independently:

\begin{description}[leftmargin=2.4cm, labelwidth=2.2cm]
  \item[Capture]        Ingestion of raw video, temporal decimation, and candidate-frame storage.
  \item[Reconstruction] Keyframe selection, \acrshort{colmap} \acrshort{sfm}, and
                        \acrshort{3dgs} model training — the photogrammetric core.
  \item[Segmentation]   Object-level decomposition of the Gaussian field via \acrshort{sam}~2
                        and cluster validation.
  \item[Generative-Reconstruction] Per-object orbit rendering, Hunyuan3D~2.0 mesh generation,
                        and \acrshort{pbr} texturing via FLUX inpainting.
  \item[Meshing]        Whole-scene surface reconstruction using one of four configurable backends.
  \item[Assembly]       Construction of the hierarchical \acrshort{usd} scene graph and
                        emplacement of reconstructed objects at \acrshort{colmap}-derived centroids.
  \item[Delivery]       Export to \texttt{.ksplat} for browser streaming, plus optional
                        \acrshort{webgl}/\acrshort{webgpu} viewer deployment.
\end{description}

Bounded contexts communicate through well-typed data contracts: frame manifests (JSON),
camera parameter files (\texttt{transforms.json}), Gaussian \texttt{.ply} archives, mask
HDF5 files, and \acrshort{usd} stage references. No context reads another's internal state directly.

\subsection{Architecture Decision Records}

All non-trivial design choices are captured as \acrfull{adr}s, versioned alongside the
source code. Two records are directly relevant to the compute fabric described in
Section~\ref{sec:containers}:

\begin{itemize}
  \item \textbf{ADR-001} (\emph{Container isolation per CUDA ABI}) — mandates one Docker
        image per incompatible CUDA/\-PyTorch combination, communicating via a shared NVMe
        volume and a lightweight REST job-handoff protocol, rather than attempting
        multi-environment conda workarounds inside a single container.
  \item \textbf{ADR-002} (\emph{GPU affinity assignment}) — binds the primary reconstruction
        workload exclusively to GPU~0 and delegates all sidecar meshing to GPU~1, preventing
        VRAM contention during the densification-intensive \acrshort{3dgs} training phase.
\end{itemize}

\subsection{Agentic Orchestration}

Conventional multi-stage photogrammetry pipelines are driven by shell scripts that invoke
tools in a fixed order and abort on the first non-zero exit code. \texttt{gaussian-toolkit}
replaces this with an \acrshort{llm} orchestrator connected to \acrshort{mcp} server
running on port 45677, which exposes over 70 tools covering every stage. The orchestrator
reasons about intermediate quality metrics — checking reprojection error before proceeding
to training, validating cluster coverage before initiating generative reconstruction — and
can reroute execution, adjust hyperparameters, or surface diagnostic messages to the
operator without human intervention. This makes the system resilient to the variability
inherent in real gallery footage: inconsistent lighting, partial occlusion, and
camera shake that would cause a brittle script to fail silently.

% ─────────────────────────────────────────────────────────────────────────────
\section{The Eight-Stage Pipeline}\label{sec:pipeline-stages}
% ─────────────────────────────────────────────────────────────────────────────

Figure~\ref{fig:pipeline-flow} illustrates the eight stages executed for a single
acquisition. Stages 1–4 constitute the photogrammetric core and are described in
detail below. Stages 5–8 are documented in the dedicated chapters that follow
(Chapter~\ref{ch:scenegraph} covers segmentation and \acrshort{usd} assembly;
Chapter~\ref{ch:meshing} covers environment meshing); a brief description of each
is included here for completeness.

\begin{figure}[htbp]
\centering
\begin{tikzpicture}[
  node distance = 0.55cm and 0.0cm,
  stage/.style  = {draw, rounded corners=4pt, minimum width=2.9cm,
                   minimum height=1.0cm, align=center, font=\small\bfseries,
                   fill=LightGrey, draw=MidGrey, text=DarkBg},
  arrow/.style  = {-Stealth, thick, color=MidGrey},
  label/.style  = {font=\scriptsize, text=MidGrey, align=center}
]
  % Row 1: stages 1-4
  \node[stage, fill=DreamLabAccent!15, draw=DreamLabAccent] (s1) {Stage 1\\Frame Extraction};
  \node[stage, fill=DreamLabAccent!15, draw=DreamLabAccent, right=of s1] (s2) {Stage 2\\Keyframe Selection};
  \node[stage, fill=DreamLabAccent!15, draw=DreamLabAccent, right=of s2] (s3) {Stage 3\\SfM (COLMAP)};
  \node[stage, fill=DreamLabAccent!15, draw=DreamLabAccent, right=of s3] (s4) {Stage 4\\3DGS Training};

  % Row 2: stages 5-8 (below, right-to-left for visual return)
  \node[stage, fill=UoSRed!10, draw=UoSRed, below=1.4cm of s4] (s5) {Stage 5\\SAM2 Segmentation};
  \node[stage, fill=UoSRed!10, draw=UoSRed, left=of s5] (s6) {Stage 6\\Gen.\ Reconstruction};
  \node[stage, fill=UoSRed!10, draw=UoSRed, left=of s6] (s7) {Stage 7\\Env.\ Meshing};
  \node[stage, fill=SuccessGreen!15, draw=SuccessGreen, left=of s7] (s8) {Stage 8\\USD Assembly};

  % Row 1 arrows
  \draw[arrow] (s1) -- (s2);
  \draw[arrow] (s2) -- (s3);
  \draw[arrow] (s3) -- (s4);

  % Turn arrow s4 -> s5
  \draw[arrow] (s4) -- (s5);

  % Row 2 arrows (right to left)
  \draw[arrow] (s5) -- (s6);
  \draw[arrow] (s6) -- (s7);
  \draw[arrow] (s7) -- (s8);

  % Annotations below row 1
  \node[label, below=0.12cm of s1] {2 fps $\to$ 600\\candidate frames};
  \node[label, below=0.12cm of s2] {Fibonacci scoring\\$\to$ 240 keyframes};
  \node[label, below=0.12cm of s3] {232/240 registered\\0.64\,px error};
  \node[label, below=0.12cm of s4] {1.18M Gaussians\\29.6\,dB PSNR};

  % Annotations above row 2
  \node[label, above=0.12cm of s5] {98.3\% coverage\\37 clusters};
  \node[label, above=0.12cm of s6] {64 orbit views\\Hunyuan3D 2.0};
  \node[label, above=0.12cm of s7] {4 backends\\TSDF/MILo/CoMe/GW};
  \node[label, above=0.12cm of s8] {31 object prims\\variant sets};

\end{tikzpicture}
\caption{The eight-stage \texttt{gaussian-toolkit} pipeline. Stages 1–4 (blue) form the
  photogrammetric core; stages 5–7 (red) perform semantic decomposition, generative
  reconstruction, and surface meshing; stage 8 (green) assembles the hierarchical
  \acrshort{usd} scene graph.}
\label{fig:pipeline-flow}
\end{figure}

\subsection{Stage 1 — Frame Extraction}

Raw 4K gallery footage is ingested via \texttt{ffmpeg} with a temporal decimation rate
of 2 frames per second. For a typical 5-minute acquisition walk this yields approximately
600 candidate frames, stored as lossless PNG at full sensor resolution. The module
\texttt{frame\_extractor.py} writes a JSON manifest alongside each batch recording
source timecode, frame index, and a SHA-256 hash of the pixel buffer. The manifest
serves as a checkpoint: if the pipeline is interrupted and resumed, frames already
extracted are skipped without re-decoding.

\subsection{Stage 2 — Keyframe Selection}

Six hundred candidate frames are far too many for efficient \acrshort{sfm} registration
and introduce redundant viewpoints that inflate the subsequent Gaussian count without
improving quality. The module \texttt{fibonacci\_sampler.py} implements a Fibonacci-sphere
viewpoint scoring heuristic in which each candidate is assigned a composite score

\[
  s_i = 0.6\,\sigma_i + 0.4\,\kappa_i
\]

\noindent where $\sigma_i$ is a normalised sharpness score (Laplacian variance, max-pooled
over $32\times32$ tiles) and $\kappa_i$ is a coverage score measuring angular distance from
the nearest already-selected keyframe on the Fibonacci sphere approximation of the view
hemisphere. The top 240 frames by composite score are retained, reducing the set by 60\%
while preserving broad angular coverage. In practice this selection consistently improves
\acrshort{sfm} registration rate compared to uniform temporal subsampling, because blurry
or redundant frames that degrade feature matching are preferentially excluded.

\subsection{Stage 3 — Structure-from-Motion}

The 240 selected keyframes are passed to \acrshort{colmap}~\citep{Schoenberger2016sfm,
Schoenberger2016mvs} running in sequential-mapper mode with SIFT~\citep{Lowe2004} features.
For the reference gallery acquisition, 232 of 240 frames registered successfully
(96.7\,\%), yielding a sparse point cloud of approximately 210,000 points with a
mean reprojection error of 0.64\,px — comfortably within the $<1$\,px threshold that
reliably produces well-conditioned Gaussian initialisation. The eight unregistered frames
were all affected by severe motion blur during camera panning at the gallery perimeter;
the Fibonacci scoring step eliminated most such frames, but eight passed the sharpness
threshold because blur was localised to a small tile.

Camera intrinsics and extrinsics are exported to \texttt{transforms.json} in the
NeRFStudio/\allowbreak{}LichtFeld-Studio format. A secondary \texttt{cameras.bin} export
is retained for downstream TSDF meshing, which requires COLMAP-native binary format.

\subsection{Stage 4 — 3DGS Training}

Gaussian splatting training is performed by \texttt{LichtFeld Studio}~\citep{LichtFeldStudio2026},
the Inria-licensed viewer and training harness that forms the upstream of the
\texttt{gaussian-toolkit} fork. Training runs for 30,000 iterations with the
\acrfull{mrnf} densification strategy, which applies a surface-normal consistency
filter during the adaptive density control loop to suppress needle-like floater Gaussians
in low-texture regions such as painted wall surfaces.

On the reference hardware (GPU~0, RTX~6000 Ada, 48\,GB VRAM) training completes in
24 minutes, converging to 1.18 million Gaussians. Hold-out evaluation on 12 reserved
frames yields:

\begin{itemize}
  \item \acrshort{psnr}: $29.6$\,dB
  \item \acrshort{ssim}: $0.921$
  \item \acrshort{lpips}: $0.108$
\end{itemize}

The trained splat renders at 160 frames per second at 1080p on the same GPU, meeting the
real-time interactive threshold with comfortable headroom for simultaneous web streaming.
The model is exported both as a raw \texttt{.ply} (for downstream segmentation and meshing)
and as a compressed \texttt{.ksplat} (for browser delivery via \acrshort{webgl}).

\subsection{Stages 5–8 — Overview}

\textbf{Stage~5 (Segmentation)} applies \acrshort{sam}~2~\citep{Ravi2024SAM2} across all
240 registered views to produce per-Gaussian cluster assignments. The full segmentation
workflow, cluster validation criteria, and the mask projection module
\texttt{mask\_projector.py} are described in Chapter~\ref{ch:scenegraph}.

\textbf{Stage~6 (Generative reconstruction)} renders 64 orbit views per retained cluster
and passes them to Hunyuan3D~2.0~\citep{Zhao2025Hunyuan} via ComfyUI (port 8188) to
generate a watertight \acrshort{pbr}-textured mesh for each object. FLUX~\citep{FLUX2024}
inpainting resolves seam artefacts. This stage is also documented in
Chapter~\ref{ch:scenegraph}.

\textbf{Stage~7 (Environment meshing)} reconstructs the room envelope and large static
surfaces using one of four selectable backends (\acrshort{tsdf}~\citep{Curless1996},
\acrshort{milo}~\citep{Guedon2025MILo}, \acrshort{come}~\citep{Radl2025CoMe}, or
Gaussian Wrapping). Backend selection and quality trade-offs are analysed in
Chapter~\ref{ch:meshing}.

\textbf{Stage~8 (USD assembly)} constructs a hierarchical \acrshort{usd} scene
graph~\citep{Pixar2019} with 31 object prims emplaced at \acrshort{colmap}-derived
centroids. Each prim carries a variant set offering both a Gaussian splat representation
and a mesh representation, enabling downstream consumers to select the appropriate fidelity
tier. Assembly is also covered in Chapter~\ref{ch:scenegraph}.

% ─────────────────────────────────────────────────────────────────────────────
\section{Three-Container Compute Fabric}\label{sec:containers}
% ─────────────────────────────────────────────────────────────────────────────

The three meshing backends — \acrshort{milo}, \acrshort{come}, and Gaussian Wrapping —
each require PyTorch extension modules compiled against specific CUDA toolkit versions
that are mutually incompatible: \acrshort{milo} requires CUDA~11.8 and PyTorch~2.0,
\acrshort{come} requires CUDA~12.1 and PyTorch~2.3, and the primary
\texttt{gaussian-toolkit} container targets CUDA~12.8 and PyTorch~2.7 for
access to the latest Hopper/Ada optimisations in the Gaussian rasteriser.
Installing all three in a single environment via conda is not feasible without
breaking at least one set of compiled extensions.

ADR-001 resolves this by mandating one Docker image per CUDA ABI. The three containers
communicate through two mechanisms: a shared NVMe Docker volume
(\texttt{gaussian-shared-vol}) that carries large binary artefacts (Gaussian \texttt{.ply}
files, depth maps, mesh outputs), and a lightweight REST job-handoff protocol in which
\texttt{gaussian-toolkit} POSTs a job descriptor to each sidecar container's local
Flask server and polls for completion. This keeps container interfaces narrow and avoids
the complexity of a full message broker for what is, in practice, a sequential handoff
with at most two concurrent sidecar jobs.

Table~\ref{tab:containers} summarises the three containers and their responsibilities.

\begin{table}[htbp]
\centering
\caption{Three-container compute fabric: CUDA environments and responsibilities.}
\label{tab:containers}
\begin{tabular}{@{}L{3.2cm} C{1.5cm} C{1.5cm} C{1.4cm} L{6.0cm}@{}}
\toprule
\textbf{Container} & \textbf{CUDA} & \textbf{Python} & \textbf{GPU} & \textbf{Responsibilities} \\
\midrule
\texttt{gaussian-toolkit} & 12.8 & 3.12 & GPU~0 &
  Orchestration, frame extraction, \acrshort{colmap} \acrshort{sfm},
  \acrshort{3dgs} training (\acrshort{mrnf} densification),
  \acrshort{tsdf} meshing, \acrshort{sam}~2 segmentation,
  mask projection, \acrshort{usd} assembly,
  \texttt{.ksplat} export, web UI and ttyd terminal. \\[4pt]
\texttt{milo} & 11.8 & 3.10 & GPU~1 &
  \acrfull{milo} mesh-in-the-loop optimisation;
  Gaussian Wrapping backend. \\[4pt]
\texttt{come} & 12.1 & 3.10 & GPU~1 &
  \acrfull{come} confidence-based surface extraction. \\
\bottomrule
\end{tabular}
\end{table}

\begin{figure}[htbp]
\centering
\begin{tikzpicture}[
  node distance  = 1.0cm and 1.2cm,
  container/.style = {draw, rounded corners=6pt, minimum width=3.6cm,
                      minimum height=2.4cm, align=center, font=\small\bfseries,
                      text=DarkBg},
  gpu/.style       = {draw, rounded corners=3pt, minimum width=3.0cm,
                      minimum height=0.9cm, align=center, font=\footnotesize,
                      fill=DarkBg!12, draw=DarkBg},
  vol/.style       = {draw, rounded corners=3pt, minimum width=3.2cm,
                      minimum height=1.0cm, align=center, font=\small,
                      fill=WarningAmber!20, draw=WarningAmber},
  arrow/.style     = {-Stealth, thick, color=MidGrey},
  restarrow/.style = {Stealth-Stealth, thick, color=DreamLabAccent, dashed}
]

  % Primary container
  \node[container, fill=DreamLabAccent!15, draw=DreamLabAccent] (gt) {
    \texttt{gaussian-toolkit}\\[2pt]
    CUDA 12.8 / Py 3.12\\
    \footnotesize Orchestration, SfM,\\
    \footnotesize 3DGS, TSDF, SAM2,\\
    \footnotesize USD, web
  };

  % Sidecar containers
  \node[container, fill=UoSRed!10, draw=UoSRed, right=2.2cm of gt] (milo) {
    \texttt{milo}\\[2pt]
    CUDA 11.8 / Py 3.10\\
    \footnotesize MILo +\\
    \footnotesize GaussianWrapping
  };

  \node[container, fill=UoSRed!10, draw=UoSRed, below=1.4cm of milo] (come) {
    \texttt{come}\\[2pt]
    CUDA 12.1 / Py 3.10\\
    \footnotesize CoMe meshing
  };

  % Shared volume
  \node[vol, below=1.8cm of gt] (vol) {
    \texttt{gaussian-shared-vol}\\
    \footnotesize NVMe shared volume\\
    \footnotesize (.ply, depth maps, meshes)
  };

  % GPU nodes
  \node[gpu, below=3.6cm of gt] (gpu0) {GPU 0 — RTX 6000 Ada\\48\,GB VRAM};
  \node[gpu, right=1.4cm of gpu0] (gpu1) {GPU 1 — RTX 6000 Ada\\48\,GB VRAM};

  % Arrows: REST job handoff
  \draw[restarrow] (gt.east) -- node[above,font=\scriptsize,text=DreamLabAccent]{REST} (milo.west);
  \draw[restarrow] (gt.east) -- ++(0.9,0) |- node[right,font=\scriptsize,text=DreamLabAccent]{REST} (come.west);

  % Arrows: shared volume
  \draw[arrow] (gt.south) -- (vol.north);
  \draw[arrow] (vol.east) -- ++(0.4,0) |- (milo.south);
  \draw[arrow] (vol.east) -- ++(0.4,0) |- (come.south);

  % GPU assignments
  \draw[arrow, color=DreamLabAccent] (gt.south) -- ++(0,-0.15) -- (gpu0.north);
  \draw[arrow, color=UoSRed] (milo.south) -- ++(0,-0.2) -- (gpu1.north);
  \draw[arrow, color=UoSRed] (come.south) -- ++(0,-0.1) -- (gpu1.north);

\end{tikzpicture}
\caption{Three-container compute fabric. The primary \texttt{gaussian-toolkit} container
  (blue) orchestrates sidecar containers \texttt{milo} and \texttt{come} (red) via REST
  job descriptors and a shared NVMe volume. GPU affinity follows ADR-002: GPU~0 for
  reconstruction, GPU~1 for sidecar meshing.}
\label{fig:containers}
\end{figure}

% ─────────────────────────────────────────────────────────────────────────────
\section{Hardware Platform}\label{sec:hardware}
% ─────────────────────────────────────────────────────────────────────────────

All experiments reported in this document were conducted on a single workstation
configured as follows:

\begin{itemize}
  \item \textbf{GPUs}: dual NVIDIA RTX~6000 Ada Generation, 48\,GB GDDR6 ECC each
        (96\,GB total \acrshort{vram}). GPU~0 handles the complete reconstruction
        pipeline; GPU~1 is reserved exclusively for sidecar meshing workloads,
        preventing the memory pressure that would otherwise arise from running
        \acrshort{milo} or \acrshort{come} concurrently with a live \acrshort{3dgs}
        model in GPU~0 VRAM.
  \item \textbf{CPU}: AMD Threadripper PRO, 48 physical cores. The high core count
        benefits the \acrshort{colmap} feature-matching and bundle-adjustment phases,
        which parallelise efficiently across CPU threads before the GPU rasteriser
        dominates in later stages.
  \item \textbf{System RAM}: 251\,GB DDR5 ECC. Large system RAM allows the full
        candidate frame set (600 frames at 4K PNG, approximately 28\,GB) to reside
        in RAM throughout the pipeline, eliminating NVMe I/O latency during repeated
        reads by the SIFT extractor.
  \item \textbf{Storage}: NVMe SSD array hosting the shared Docker volume; sequential
        read bandwidth in excess of 12\,GB/s sustains the depth-map streaming rates
        required by \acrshort{tsdf} fusion without introducing pipeline stalls.
\end{itemize}

The 96\,GB combined \acrshort{vram} budget is not exhausted by the reference acquisition:
\acrshort{3dgs} training with 1.18\,M Gaussians peaks at approximately 18\,GB on GPU~0,
leaving substantial headroom for larger scenes or higher Gaussian counts in future
acquisitions. GPU~1 headroom (typically 6–12\,GB utilised by sidecar meshing) permits
the concurrent execution of multiple meshing backend comparisons without swapping.

% ─────────────────────────────────────────────────────────────────────────────
\section{Agentic Control Surface}\label{sec:control-surface}
% ─────────────────────────────────────────────────────────────────────────────

The operator's primary interface to the running pipeline is the \acrshort{mcp} server
embedded in \texttt{gaussian-toolkit}, which listens on port~45677 and exposes
over 70 typed tool calls. These tools span every bounded context: from
\texttt{ingest\_video} and \texttt{run\_colmap} at the front of the pipeline,
through \texttt{train\_3dgs}, \texttt{run\_sam2}, and \texttt{dispatch\_hunyuan},
to \texttt{assemble\_usd} and \texttt{export\_ksplat} at the delivery end.

The \acrshort{llm} orchestrator connects to this \acrshort{mcp} server and drives
the pipeline by issuing tool calls, inspecting return values (which include structured
quality metrics and error codes), and deciding whether to proceed, adjust a
hyperparameter, or surface an alert. This architecture cleanly separates the
\emph{reasoning layer} (the orchestrator, which can be swapped for any
\acrshort{mcp}-compatible client) from the \emph{execution layer} (the toolkit,
which exposes a stable tool API). Operators can also invoke tools directly from any
\acrshort{mcp} client, enabling ad-hoc inspection and intervention without modifying
the pipeline code.

Two secondary control surfaces are provided for human operators who prefer a graphical
or terminal interface. The web \acrshort{ui} (Figure~\ref{fig:webui}) presents a
dashboard summarising pipeline state, live training metrics (loss, PSNR, iteration
count), and GPU utilisation. The ttyd terminal (Figure~\ref{fig:ttyd}) exposes a full
shell inside the \texttt{gaussian-toolkit} container through a browser tab, enabling
direct inspection of intermediate artefacts, log files, and the \acrshort{colmap}
workspace without requiring SSH access.

\begin{figure}[htbp]
  \centering
  \includegraphics[width=0.8\textwidth]{v2/webui.png}
  \caption{The \texttt{gaussian-toolkit} web control interface, showing pipeline stage
    status, live \acrshort{3dgs} training metrics, and \acrshort{gpu} telemetry.}
  \label{fig:webui}
\end{figure}

\begin{figure}[htbp]
  \centering
  \includegraphics[width=0.8\textwidth]{v2/ttyd.png}
  \caption{The ttyd browser terminal embedded in the \texttt{gaussian-toolkit} web
    interface, providing in-browser shell access to the container for direct artefact
    inspection and ad-hoc tool invocation.}
  \label{fig:ttyd}
\end{figure}

\begin{keyfinding}
  The combination of \acrshort{mcp}-mediated agentic orchestration with direct web and
  terminal control surfaces allows \texttt{gaussian-toolkit} to be operated both
  fully autonomously (the orchestrator drives all 8 stages without human input) and
  interactively (an operator inspects intermediate artefacts and overrides decisions at
  any stage boundary). Neither mode requires modifying pipeline code.
\end{keyfinding}
